\documentclass[11pt]{article}

\usepackage{amsmath}
\usepackage{amssymb}
\usepackage[margin=1in]{geometry}
\usepackage{booktabs}
\usepackage{hyperref}

\title{Reference frames and radiation moments in the Monte Carlo module}
\author{Athena++ Monte Carlo radiation transport}
\date{}

\begin{document}
\maketitle

\begin{abstract}
This note records the frame transformations used by the Monte Carlo radiation transport
module and the estimator used to accumulate radiation moments. It is written to be read
alongside the source. The material is split across three files by layer:
\texttt{src/monte\_carlo/tetrad.cpp} holds the geometric primitives of
Sec.~\ref{sec:construct}, acting on bare four-vectors and a metric;
\texttt{src/monte\_carlo/photon\_frames.cpp} applies them to photons and to accumulated
moments, as in Secs.~\ref{sec:project} and \ref{sec:transform}; and the estimator of
Sec.~\ref{sec:accumulate} is assembled in the \texttt{UpdateMoments} routine of
\texttt{src/monte\_carlo/montecarloblock.cpp}. The emphasis is on which quantity is
expressed in which frame, since almost every defect found in this code has come from a
frame convention being reimplemented in a second place and drifting from the first.
\end{abstract}

\section{Conventions}

The metric signature is $(-,+,+,+)$. Greek indices $\mu,\nu$ run over coordinate
components and Latin indices in parentheses, $(a),(b)$, over orthonormal tetrad
components. A photon four-momentum is $k^\mu$; the transport is null, so
\begin{equation}
  g_{\mu\nu} k^\mu k^\nu = 0 ,
  \label{eq:null}
\end{equation}
and this is satisfied to machine precision at the photon position throughout a run.

Monte Carlo quantities are in cgs; the hydrodynamic grid is in code units, and
\texttt{l\_cgs} converts a coordinate length to centimetres. Photon weights $w$ are
physical photon numbers per sample, not probabilities.

\subsection{Photon storage}

The stored time component \texttt{k0p} always holds the photon energy in whatever frame
the photon is currently expressed in; \texttt{ep} is an alias for the same storage. The
spatial components are stored one of two ways. The general pusher keeps dimensional
coordinate components $k^i$. The legacy pushers, and every frame after a transformation
to the comoving frame, keep a unit propagation direction $n^i$. The accessors
\texttt{Photon::GetFourVector} and \texttt{SetFourVector} hide the difference by
reconstructing
\begin{equation}
  k^\mu = \left( E,\; k^i \right)
  \qquad\text{or}\qquad
  k^\mu = \left( E,\; E\, n^i \right) ,
\end{equation}
so that everything downstream sees a genuine contravariant four-vector.

\subsection{Path length as handed to the moments}

The pushers advance a photon and then report the segment length. The legacy pushers
integrate a coordinate path length directly and pass \texttt{dl * l\_cgs}. The general
pusher integrates the geodesic in an affine parameter $\lambda$, with
$dx^\mu = k^\mu \, d\lambda$, and passes \texttt{step * ep * l\_cgs}. Since the time
component of that relation is $dt = k^0 d\lambda = E\, d\lambda$, both pushers deliver the
same thing: a path length measured in the \emph{coordinate} frame, in centimetres. Call it
$\Delta \ell_{\rm coord}$. Every frame-dependent rescaling below starts from it.

\section{The three frames}

\begin{table}[h]
\centering
\begin{tabular}{lll}
\toprule
Frame & Orthonormal & Definition \\
\midrule
\texttt{MCFRAME\_LAB}      & yes & normal (Eulerian) observer in GR; Eulerian frame in flat spacetime \\
\texttt{MCFRAME\_COMOVING} & yes & rest frame of the four-velocity held in \texttt{vel} \\
\texttt{MCFRAME\_COORD}    & no  & the raw coordinate basis \\
\bottomrule
\end{tabular}
\end{table}

The coordinate basis is retained because the components of $T^{\mu\nu}$ in it are what
couple to the GRMHD conservative variables and the radiation four-force. Those components
are dimensionally inhomogeneous by construction, since $k^\theta$ and $k^\phi$ carry
inverse length, and no attempt is made to normalise them.

\section{Frame construction}
\label{sec:construct}

\subsection{General relativity}

Both orthonormal frames are built by Gram--Schmidt orthonormalisation against a timelike
four-velocity. Given $u^\mu$, the timelike leg is
\begin{equation}
  e_{(0)}^{\ \ \mu} = \frac{u^\mu}{\sqrt{\left| g_{\alpha\beta} u^\alpha u^\beta \right|}} ,
\end{equation}
and the three spacelike legs are obtained by orthonormalising the coordinate directions
$\delta^\mu_{\ 1}, \delta^\mu_{\ 2}, \delta^\mu_{\ 3}$ against $e_{(0)}$ and against each
other. The covariant legs $e^{(a)}_{\ \ \mu}$ then project a coordinate four-vector onto
orthonormal components,
\begin{equation}
  p^{(a)} = e^{(a)}_{\ \ \mu} \, k^\mu .
  \label{eq:project}
\end{equation}

The lab frame uses the normal observer,
\begin{equation}
  n^\mu = -\alpha \, g^{\mu t} , \qquad \alpha = \frac{1}{\sqrt{-g^{tt}}} ,
\end{equation}
which stays well defined inside the ergosphere. The comoving frame uses the four-velocity
stored in \texttt{vel}: the fluid when boosts are enabled, and the normal observer
otherwise, in which case the two frames coincide by construction.

Both are evaluated once per zone, at the zone centre, in
\texttt{ComputeTransformations()}, and stored as the arrays \texttt{boost\_lab} and
\texttt{boost\_cmv}. A single matrix multiply then carries a coordinate four-vector to
orthonormal components.

It is worth stating what this replaced, because the failure was not obvious. Earlier
versions built \texttt{boost\_cmv} as a flat Lorentz boost from the components of
\texttt{vel}, treating $u^0$ as $\gamma$ and $u^i$ as $\gamma \beta^i$. In general
relativity \texttt{vel} holds a coordinate-frame four-velocity, so those components are
not orthonormal: for a static observer in Kerr--Schild the spatial part carries the shift
vector and the components miss the Minkowski norm by $25\%$ at $r=6M$. The construction
was not a Lorentz transformation at all, and a fluid at rest acquired a spurious Doppler
factor.

\subsection{Flat spacetime}

In flat spacetime the same two arrays hold a closed-form Lorentz boost built from
$u^\mu = \gamma(1, \beta^i)$,
\begin{equation}
  \Lambda^{(0)}_{\ \ (0)} = \gamma , \qquad
  \Lambda^{(0)}_{\ \ (i)} = -\gamma \beta_i , \qquad
  \Lambda^{(i)}_{\ \ (j)} = \delta_{ij} + \frac{\gamma^2 \beta_i \beta_j}{1+\gamma} ,
\end{equation}
which acts on the lab-frame components rather than on coordinate components. This is the
special case of the tetrad construction above with $g_{\mu\nu} = \eta_{\mu\nu}$, and the
two agree to machine precision where both apply.

\section{Projecting a photon}
\label{sec:project}

For a photon in a given zone, \texttt{PhotonFrames} returns three numbers per frame: an
energy $E$, a unit propagation direction $n^i$, and a path length $\Delta \ell$. The
frames are filled lazily and cached, so several consumers asking for the same frame pay
for one projection.

\subsection{Energy and direction}

In the orthonormal frames the projection is Eq.~\eqref{eq:project}, followed by
\begin{equation}
  E = p^{(0)} , \qquad
  n^i = \frac{p^{(i)}}{\sqrt{\delta_{jk}\, p^{(j)} p^{(k)}}} .
  \label{eq:normalise}
\end{equation}

The normalisation in Eq.~\eqref{eq:normalise} deserves comment, because the more obvious
choice $n^i = p^{(i)}/p^{(0)}$ is wrong here. If the projection were exactly null then the
two agree, since $\eta_{ab} p^{(a)} p^{(b)} = 0$ implies
$\delta_{jk} p^{(j)} p^{(k)} = (p^{(0)})^2$. The photon four-vector does satisfy
Eq.~\eqref{eq:null} at its own position, but the tetrad is orthonormal with respect to the
metric at the \emph{zone centre}, so the projected vector is null only to the accuracy with
which the metric is constant across a zone. Measured on the isothermal shell problem, the
null violation at the photon is $1.2\times 10^{-16}$ while the same contraction against the
zone-centre metric is $1.1\times 10^{-1}$, falling roughly first order under angular
refinement. Dividing by $p^{(0)}$ would therefore leave $n^i$ off unit length by that
amount, and the resulting moment tensor would not be traceless. Dividing by the spatial
magnitude confines the error to the direction, where it belongs, and preserves
Eq.~\eqref{eq:traceless} below exactly.

In the coordinate frame no normalisation is possible or wanted, and the stored direction
is simply $k^i/k^0$.

\subsection{Path length}

Because $dx^\mu = k^\mu d\lambda$, the interval measured in a frame whose timelike leg
gives the photon energy $E$ is $\Delta \ell = E \, \Delta\lambda$. The affine step is
frame independent, so the ratio between any two frames is the ratio of the energies. With
$E_{\rm coord} = k^0$ the stored energy in the coordinate frame,
\begin{equation}
  \Delta \ell = \Delta \ell_{\rm coord} \, \frac{E}{E_{\rm coord}} .
  \label{eq:dl}
\end{equation}
This factor is easy to omit, and omitting it is not visible in flat spacetime, where the
lab tetrad is trivial and $E = E_{\rm coord}$. It is the reason the path length is carried
in the same structure as the energy rather than passed alongside it: handing a consumer a
frame-correct energy next to a coordinate path length simply relocates the trap.

\section{Moment accumulation}
\label{sec:accumulate}

The radiation stress tensor is accumulated with the covariant estimator
\begin{equation}
  T^{(a)(b)} = \frac{1}{\Delta V_4} \sum_{\rm photons} w \; p^{(a)} p^{(b)} \, \Delta\lambda ,
  \label{eq:estimator}
\end{equation}
where $\Delta V_4$ is the four-volume of the zone over the integration window and the sum
runs over every segment of every photon track crossing it. Writing
$p^{(a)} = E\,(1, n^i)$ and using $\Delta\lambda = \Delta \ell / E$, the per-segment
contribution reduces to a single weight,
\begin{equation}
  \text{weight} = \frac{w \, E \, \Delta \ell}{c} ,
\end{equation}
multiplied into the products of the direction components. The three stored quantities are
\begin{equation}
  E_{\rm r} = T^{(0)(0)} , \qquad
  F_{\rm r}^{i} = c \, T^{(0)(i)} , \qquad
  P_{\rm r}^{ij} = T^{(i)(j)} ,
\end{equation}
so that in code they become
\begin{equation}
  E_{\rm r} \mathrel{+}= \text{weight} , \qquad
  F_{\rm r}^{i} \mathrel{+}= \text{weight} \; n^i \, c , \qquad
  P_{\rm r}^{ij} \mathrel{+}= \text{weight} \; n^i n^j .
\end{equation}
The mapping from $(a,b)$ to a storage slot, and which slots carry the extra factor of $c$,
is held in a single table so that the accumulation and the frame transform of
Sec.~\ref{sec:transform}, which are inverses of each other, cannot disagree.

Because the spatial direction is a unit vector, the stress tensor of a null-particle gas is
traceless in any orthonormal frame,
\begin{equation}
  \eta_{ab} T^{(a)(b)} = -E_{\rm r} + P_{\rm r}^{11} + P_{\rm r}^{22} + P_{\rm r}^{33} = 0 ,
  \label{eq:traceless}
\end{equation}
which is checked directly by the frame-consistency tests and is a sharp diagnostic in
curved spacetime.

\section{Transforming accumulated moments}
\label{sec:transform}

The moments are a rank-two tensor and the transformation between two frames is the same
matrix for every photon in a zone. The sum and the transformation therefore commute,
\begin{equation}
  \sum_{\rm photons} \Lambda p \, \Lambda p
  \;=\; \Lambda \left( \sum_{\rm photons} p \, p \right) \Lambda^{\mathsf T} ,
\end{equation}
so the estimator of Eq.~\eqref{eq:estimator} can be evaluated once per zone at output
rather than accumulated photon by photon,
\begin{equation}
  T^{(a')(b')}_{\rm com} = \Lambda^{(a')}_{\ \ (a)} \Lambda^{(b')}_{\ \ (b)} \, T^{(a)(b)}_{\rm lab} .
\end{equation}
This is exact rather than an approximation of the accumulation, and it removes a second
per-photon projection.

The transformation is obtained by composing the two tetrads,
\begin{equation}
  \Lambda^{(a')}_{\ \ (a)} = e^{(a')}_{{\rm fluid}\ \mu} \; e_{{\rm lab}\ (a)}^{\ \ \ \ \mu} ,
\end{equation}
and \emph{not} by rebuilding a pure boost from a three-velocity. Both frames are
orthonormal at the same event, so the map is certainly a Lorentz transformation, and it
satisfies $\Lambda^{\mathsf T} \eta \, \Lambda = \eta$ to machine precision. It is not in
general a pure boost. The two spatial triads are grown by Gram--Schmidt from the same
coordinate trial vectors but against different timelike legs, so they differ by a rotation
as well as a boost whenever the relative velocity has all three spatial components; the
antisymmetric part of the spatial block is $6.9\%$ for a fully three-dimensional velocity
in Kerr--Schild, and vanishes identically when the velocity is aligned with a coordinate
axis. Using a pure boost would leave $E_{\rm r}$ correct, since it is a scalar under that
rotation, while silently rotating the axes of $F_{\rm r}^i$ and $P_{\rm r}^{ij}$.

\section{Normalisation}

Photon weights are set from a number emissivity, which is naturally a comoving-frame
quantity, multiplied by the cell volume and the integration time,
\begin{equation}
  w = \eta_{\rm com} \; V \; \Delta t .
\end{equation}
Mixing a comoving emissivity with a coordinate volume and a coordinate time is correct
rather than approximate, because the four-volume element is invariant,
\begin{equation}
  dV_{\rm proper} \, d\tau = \sqrt{-g} \, d^3x \, dt ,
\end{equation}
so the product is the invariant photon number. In flat spacetime the same identity holds
for a different reason: the Lorentz contraction of the volume cancels the time dilation.
The cell volume returned by the coordinate class integrates $\sqrt{-g}$ exactly, so this
holds term by term in Kerr--Schild.

\texttt{NormalizeMoments} divides by the same $\Delta t \, V$ that appears in the weight.
Both factors therefore cancel, and the units of the integration time never reach the
answer; the physical rate comes entirely from the emissivity being in cgs.

\section{Approximations and their signatures}

\begin{enumerate}
\item \textbf{Zone-centred frames.} The tetrads and the boost are evaluated once per zone
at the zone centre. For the boost this is exact given the existing assumption that the
fluid velocity is uniform within a cell. For the tetrad it is a further approximation with
no counterpart in that assumption: the orthonormal basis rotates across a cell in any
curvilinear coordinate system, with or without gravity, by an angle of order the cell's
angular size. In Cartesian coordinates it is exact.

\item \textbf{Non-null projection.} A direct consequence of the above, quantified in
Sec.~\ref{sec:project} and handled by the normalisation of Eq.~\eqref{eq:normalise}.

\item \textbf{Accumulated versus derived comoving moments.} Both apply the rescaling of
Eq.~\eqref{eq:dl}. They agree exactly in flat spacetime and are checked there. In curved spacetime they differ, because the accumulated
path renormalises each photon direction in the comoving frame while a tensor transform
cannot reproduce a per-photon renormalisation. On the isothermal shell problem the
difference is $12$--$32\%$ in $E_{\rm r}$ and the pressure diagonal, and larger in the
flux, which is a near-cancelling residual. The derived form is preferred: it is an exact
transform of the lab moments and inherits their tracelessness, so it satisfies both
Eq.~\eqref{eq:traceless} and the tensor relation, whereas the accumulated form satisfies
only the first. The difference between them is a direct measure of the zone-centring error
and is otherwise difficult to observe.

\item \textbf{Polarization.} Not yet transformed between frames in the general
relativistic path. The polarization tensor is parallel transported along geodesics and has
a rank-two transformation available, but the comoving transforms do not currently apply
it.
\end{enumerate}

\end{document}
